\section{Videoquellen und Spurenbildung}
\label{sec:videoquellen}

Die Identifikation der technischen Quelle eines Videos erfordert eine Erfassung der Verarbeitungsschritte. Im forensischen Kontext hinterlässt jede Verarbeitungsstufe spezifische Spuren in der Containerstruktur, die als struktureller Fingerabdruck dienen kann \cite[S. s68]{gloe_forensic_2014}. In der Praxis lassen sich MP4-Dateien primär in native Aufnahmen, Social-Media-Verarbeitungen und synthetische Inhalte klassifizieren.

Native Videoaufnahmen aus Smartphones oder Digitalkameras weisen durch herstellerspezifische Firmware und proprietäre Metadaten (z.\,B. \(udta\)-Boxen) eine hohe strukturelle Diversität auf. Diese Varianz in der \(moov\)-Topologie ermöglicht oftmals eine präzise Zuordnung zu Gerätemodellen \cite{shullani_vision_2017, iuliani_video_2019}. Beim Transfer über Social-Media-Plattformen geht diese native Integrität durch automatische Verarbeitungen in der Regel verloren. Herstellerspezifische Boxen werden eliminiert und durch eine stark uniformierte Baumstruktur ersetzt, die wiederum charakteristisch für die jeweilige Distributionsplattform oder den eingesetzten Software-Encoder (z.\,B. \(ffmpeg\)) ist \cite[S. 947]{yang_efficient_2020}.

Generative KI-Modelle erzeugen Videos direkt cloudbasiert als MP4-Container, ohne dass physikalische Kamerahardware involviert ist \cite{yang_comprehensive_2025}. Zur Umsetzung regulatorischer Vorgaben zur Inhaltsauthentizität betten viele Generatoren wie OpenAI Sora oder Adobe Firefly kryptografische Metadaten-Manifeste nach dem C2PA-Standard ein \cite[Kap. 9]{c2pa_content_2024}. Diese umfassen vollständige Zertifikatsketten sowie digitale Signaturen, was zu strukturell massiven \(uuid\)-Boxen im Megabyte-Bereich führt. Die Platzierung dieser Validierungs-Strukturen unmittelbar am Dateianfang definiert ein distinktes syntaktisches Profil, das synthetische Container eindeutig von nativen oder netzwerk-komprimierten Videos abgrenzen kann.